\chapter{Fiscalização de Contratos de TI}

\section{Contratações de TI como objeto de controle}

Contratações de tecnologia concentram riscos de especificação inadequada, dependência do fornecedor, superdimensionamento, subutilização, medição subjetiva, segurança, proteção de dados, continuidade e sobrepreço. O auditor deve verificar não apenas se o rito formal foi cumprido, mas se a contratação resolve necessidade pública, possui critérios verificáveis e produz valor compatível com custo e risco.

\begin{teoria}[cadeia lógica da contratação]
Uma contratação tecnicamente defensável deve formar uma cadeia rastreável:
\[
\text{necessidade}\rightarrow\text{alternativas}\rightarrow\text{solução}\rightarrow\text{requisitos}\rightarrow\text{estimativa}\rightarrow\text{seleção}\rightarrow\text{execução}\rightarrow\text{medição}\rightarrow\text{pagamento}\rightarrow\text{benefício}.
\]
O auditor procura rupturas nessa cadeia. Se o quantitativo não decorre da demanda, se o preço não deriva de fontes comparáveis ou se o pagamento não decorre de entrega comprovada, existe risco de economicidade, legitimidade ou responsabilização.
\end{teoria}

\section{Lei nº 14.133/2021 aplicada à TIC}

A Lei nº 14.133/2021 reforça planejamento, governança, segregação de funções, motivação, julgamento objetivo, competitividade, economicidade, eficiência e gestão de riscos. Em TI, esses princípios precisam ser traduzidos em decisões técnicas e evidências.

\subsection{Planejamento não é formulário}

Planejar significa compreender o problema antes de escolher a tecnologia. Uma boa instrução processual deve permitir responder:
\begin{itemize}
\item qual problema público ou operacional será resolvido;
\item qual demanda histórica e futura sustenta o quantitativo;
\item quais alternativas foram comparadas e por que foram descartadas;
\item quais riscos tecnológicos, contratuais e de transição existem;
\item qual benefício esperado justifica o custo total.
\end{itemize}

Quando a solução já está definida e o ETP é escrito apenas para defendê-la, há risco de racionalização posterior e direcionamento.

\subsection{Segregação de funções}

Segregação não significa multiplicar agentes artificialmente. Significa impedir concentração incompatível de poderes de especificar, selecionar, executar, medir e pagar sem revisão proporcional ao risco. Em estruturas pequenas, controles compensatórios, supervisão e rastreabilidade podem ser necessários.

\section{IN SGD/ME nº 94/2022}

A IN SGD/ME nº 94/2022 organiza o processo de contratação de soluções de TIC no SISP em três grandes fases:

\begin{center}
\begin{tikzpicture}[node distance=7mm]
\tikzset{f/.style={draw=azulPetroleo,fill=azulClaro,rounded corners,minimum width=3.8cm,minimum height=1cm,align=center}}
\node[f] (p) {Planejamento\\da Contratação};
\node[f,right=of p] (s) {Seleção\\do Fornecedor};
\node[f,right=of s] (g) {Gestão\\do Contrato};
\draw[-Latex,thick,ambar] (p)--(s)--(g);
\node[below=7mm of s,draw=ambar,fill=ambarClaro,rounded corners,align=center] {Gerenciamento de riscos acompanha todo o ciclo};
\end{tikzpicture}
\end{center}

\subsection{Equipe de Planejamento}

A equipe deve combinar competências para compreender demanda, solução e aspectos administrativos. O auditor deve avaliar participação efetiva, registros de decisão, competência técnica e conflitos de interesse — não apenas a portaria de designação.

\subsection{Estudo Técnico Preliminar}

O ETP demonstra a necessidade e a viabilidade da solução. Em TIC, deve considerar, conforme o caso:
\begin{itemize}
\item volume de demanda, sazonalidade e crescimento;
\item alternativas on-premises, nuvem, serviço gerenciado, desenvolvimento ou aquisição;
\item interoperabilidade e integração com legado;
\item disponibilidade, RTO e RPO;
\item segurança e proteção de dados;
\item licenciamento, suporte e custo total de propriedade;
\item portabilidade e estratégia de saída;
\item capacidade interna para operar e fiscalizar.
\end{itemize}

\begin{cebraspe}
ETP e Termo de Referência não são equivalentes. O ETP justifica a solução; o TR transforma a solução escolhida em requisitos, critérios de execução, medição e pagamento.
\end{cebraspe}

\subsection{Termo de Referência}

O TR deve permitir que fornecedor, gestor, fiscal e auditor entendam exatamente o que significa ``cumprir o contrato''. Isso exige requisitos verificáveis, critérios de aceite, SLAs, responsabilidades, segurança, transição, sanções e forma de pagamento.

Um TR que especifica pessoas, mas não resultado, tende a transferir risco de produtividade para a Administração. Um TR que especifica resultado sem definir como medi-lo cria disputa de interpretação.

\section{Dimensionamento e memória de cálculo}

Quantitativos devem decorrer de premissas verificáveis. Em TI, exemplos comuns:

\begin{table}[ht]
\centering
\small
\begin{tabularx}{\textwidth}{l X X}
\toprule
Objeto & Dados relevantes & Risco de erro \\
\midrule
Licenças & usuários ativos, simultaneidade, perfis, crescimento & ociosidade ou falta de acesso. \\
Storage & consumo histórico, retenção, crescimento, compressão & capacidade ociosa cara ou insuficiente. \\
Service desk & volume, severidade, sazonalidade, automação & postos superdimensionados ou SLA inviável. \\
Cloud & consumo por serviço, egress, logs, suporte & teto subestimado e desperdício. \\
SOC & fontes, EPS/GB-dia, casos de uso, MTTD/MTTR & contratar pessoas sem cobertura real. \\
\bottomrule
\end{tabularx}
\end{table}

\begin{exemplo}[dimensionamento de licenças]
Se há 1.650 usuários ativos mensais e pico de 1.120 simultâneos, contratar 5.000 licenças exige justificar crescimento, regras do fornecedor, perfis e margens. ``Crescimento futuro'' sem série histórica não é memória de cálculo.
\end{exemplo}

\section{Pesquisa de preços — IN SEGES/ME nº 65/2021}

Pesquisa de preços é análise de mercado, não coleta mecânica de três propostas. Fontes precisam ser comparáveis e criticamente avaliadas.

\subsection{Comparabilidade antes da estatística}

Antes de calcular média ou mediana, normalize:
\begin{itemize}
\item objeto e escopo;
\item unidade e quantidade;
\item prazo de licença/assinatura;
\item suporte e manutenção;
\item SLA e requisitos de segurança;
\item data de referência e atualização;
\item impostos e condições comerciais.
\end{itemize}

A mediana é robusta a outliers, mas não corrige preços de objetos diferentes. Essa é pegadinha recorrente.

\subsection{Outliers e memória de cálculo}

Valores inexequíveis, inconsistentes ou excessivamente elevados podem demandar exclusão, desde que haja justificativa. O objetivo é que um terceiro consiga reproduzir a estimativa e entender por que determinadas fontes foram usadas ou descartadas.

\section{Competitividade e direcionamento}

Indícios que merecem teste adicional:
\begin{itemize}
\item requisitos copiados de catálogo de fabricante;
\item certificações raras sem relação com o objeto;
\item prazo de implantação compatível apenas com fornecedor incumbente;
\item prova de conceito desenhada sobre ambiente já dominado por um fornecedor;
\item agrupamento de itens independentes;
\item marca ou modelo sem motivação técnica suficiente.
\end{itemize}

Indício não é prova. O auditor deve testar necessidade, proporcionalidade e efeito sobre a competição.

\section{Gestão e fiscalização da execução}

\subsection{Ordem de Serviço e rastreabilidade}

A demanda deve ser identificável: escopo, prioridade, prazo, critérios de aceite, responsável e preço/quantidade quando aplicável. A cadeia esperada é:
\[
\text{demanda autorizada}\rightarrow\text{execução}\rightarrow\text{evidência}\rightarrow\text{aceite}\rightarrow\text{medição}\rightarrow\text{pagamento}.
\]

Se a nota fiscal não pode ser reconciliada com ordens de serviço e entregas, a fiscalização perde capacidade de demonstrar a liquidação da despesa.

\subsection{Recebimento}

Receber não é apenas assinar. Em software, aceite pode envolver funcionalidade, desempenho, segurança e documentação. Em infraestrutura, inventário, serial, configuração, garantia e teste. Em serviço, chamados, logs, métricas, artefatos e resultado.

\subsection{Evidência produzida pela contratada}

Não é automaticamente inválida. O auditor deve avaliar:
\begin{itemize}
\item possibilidade de alteração;
\item completude;
\item controle de acesso;
\item fonte primária;
\item capacidade de recálculo/corroborar;
\item conflito de interesse.
\end{itemize}

Uma planilha unilateral é fraca; logs brutos protegidos e recalculáveis são evidência mais robusta.

\section{SLAs, SLIs e desenho de incentivos}

A disponibilidade pode ser calculada por:
\[
A=\frac{T_{janela}-T_{indisponível}}{T_{janela}}\times100\%.
\]

Mas a fórmula sozinha não basta. É preciso definir janela, relógio, severidade, manutenção programada, eventos de terceiros, arredondamento e falhas do próprio monitoramento.

\begin{exemplo}[99,9\%]
Em 30 dias 24x7, há 43.200 minutos. Uma meta de 99,9\% admite, teoricamente, 43,2 minutos de indisponibilidade não excluída.
\end{exemplo}

\subsection{Gaming de indicadores}

Indicadores mudam comportamento. Medir apenas ``primeira resposta'' incentiva respostas automáticas rápidas e restauração lenta. Medir apenas ``chamados fechados'' pode incentivar fechamento prematuro e reabertura.

Um bom conjunto combina, conforme o serviço:
\begin{itemize}
\item tempo de resposta;
\item tempo de restauração/resolução;
\item reabertura;
\item backlog envelhecido;
\item reincidência;
\item satisfação/aceite;
\item impacto por severidade.
\end{itemize}

\section{Pagamento e economicidade}

Pagamento deve refletir entrega e resultado. Menor preço não prova economicidade. Uma licença barata e ociosa continua sendo desperdício; um SaaS barato no primeiro ano pode ser caro ao considerar egress, suporte e saída.

\begin{teoria}[economicidade de ciclo de vida]
Avalie:
\[
\text{custo de aquisição} + \text{operação} + \text{suporte} + \text{mudança} + \text{saída} - \text{benefícios demonstrados}.
\]
A fórmula é conceitual: o ponto é impedir análise baseada apenas no preço inicial.
\end{teoria}

\section{Alterações, prorrogação e reequilíbrio}

Alterações precisam de fundamento e motivação. Muitas mudanças podem revelar planejamento deficiente.

Prorrogação deve ser reavaliada à luz de desempenho, vantajosidade, demanda e risco de transição. Se a renovação ocorre apenas porque migrar é difícil, o contrato pode ter produzido lock-in.

Reequilíbrio econômico-financeiro exige analisar evento, alocação de riscos e impacto demonstrado. Não existe direito automático por aumento de custos ordinários.

\section{Lock-in e estratégia de saída}

Fontes comuns de lock-in:
\begin{itemize}
\item formatos proprietários;
\item APIs exclusivas;
\item dados sem exportação completa;
\item código/documentação indisponíveis;
\item conhecimento concentrado na contratada;
\item custos elevados de egress/migração;
\item chaves criptográficas controladas exclusivamente pelo fornecedor.
\end{itemize}

O plano de saída deve prever dados, metadados, código/licenças, contas, chaves, documentação, cronograma, suporte à transição, testes e eliminação segura.

\section{Segurança, LGPD e continuidade}

Contratos de TI devem incorporar riscos de segurança e proteção de dados no próprio objeto. Dependendo do serviço, podem exigir:
\begin{itemize}
\item papéis e instruções de tratamento;
\item suboperadores e transferências;
\item notificação de incidentes;
\item IAM/PAM e contas nominativas;
\item criptografia e gestão de chaves;
\item RPO/RTO e testes de restauração;
\item retenção e eliminação de dados;
\item direito de auditoria e preservação de evidência.
\end{itemize}

\begin{atencao}
Backup contratado sem restore test não demonstra recuperabilidade. RTO em cláusula sem arquitetura e teste é apenas intenção.
\end{atencao}

\section{Matriz de achados típicos}

\begin{table}[ht]
\centering
\small
\begin{tabularx}{\textwidth}{X X X}
\toprule
Condição & Risco & Evidência adicional \\
\midrule
Licenças ociosas & superdimensionamento & usuários ativos, perfis, tendência. \\
SLA unilateral & pagamento indevido & logs brutos, recálculo, integridade. \\
Sem plano de saída & lock-in & exportação, APIs, teste de migração. \\
Aceite genérico & entrega não comprovada & critérios, testes, artefatos. \\
Backup sem restore & RTO/RPO não demonstrados & restore test, tempos, logs. \\
Recursos cloud sem owner & desperdício & tags, fatura, utilização, autorização. \\
\bottomrule
\end{tabularx}
\end{table}

\section{Mapa mental funcional}

\mapafuncional{Contratos de TI}
{Necessidade pública deve permanecer rastreável até o pagamento e o benefício}
{Necessidade $\rightarrow$ ETP $\rightarrow$ TR $\rightarrow$ preço $\rightarrow$ seleção $\rightarrow$ OS $\rightarrow$ evidência $\rightarrow$ aceite $\rightarrow$ pagamento.\\Riscos atravessam todas as fases.}
{Se não há memória de cálculo, questione quantitativo.\\Se não há fonte reproduzível, questione preço/SLA.\\Se não há critério de aceite, pagamento perde base técnica.\\Se saída não foi testada, pense em lock-in.}
{Três propostas $\neq$ pesquisa adequada.\\Menor preço $\neq$ economicidade.\\SLA contratual $\neq$ capacidade demonstrada.\\Ausência de penalidade $\neq$ cumprimento.\\Backup existente $\neq$ restore comprovado.}
{Consigo refazer o quantitativo?\\Consigo refazer o preço?\\Consigo ligar nota fiscal a entrega?\\Consigo recalcular SLA?\\Consigo migrar do fornecedor?\\Consigo provar RTO/RPO?}

\begin{resumo}
\textbf{Regra de ouro para prova e auditoria:} não confunda documento com controle, cláusula com capacidade, relatório com evidência, preço baixo com economicidade ou fornecedor atual com solução inevitável.
\end{resumo}

\section{Banco de questões — Fiscalização de Contratos de TI}

\subsection*{N1 — Fundamentos}
\begin{questao}{\nivel{N1} 15.01}Segundo a IN SGD/ME nº 94/2022, o processo de contratação de TIC organiza-se em:\begin{enumerate}[label=\Alph*)]\item Planejamento da Contratação, Seleção do Fornecedor e Gestão do Contrato\item ETP, pagamento e auditoria\item orçamento, empenho e liquidação apenas\item desenvolvimento, teste e produção\item demanda, homologação e encerramento\end{enumerate}\end{questao}
\begin{questao}{\nivel{N1} 15.02}O Estudo Técnico Preliminar tem como função principal:\begin{enumerate}[label=\Alph*)]\item caracterizar a necessidade, avaliar alternativas e demonstrar viabilidade da solução\item substituir o contrato\item atestar pagamento\item aplicar penalidade\item executar aceite definitivo\end{enumerate}\end{questao}
\begin{questao}{\nivel{N1} 15.03}O Termo de Referência deve definir, entre outros elementos:\begin{enumerate}[label=\Alph*)]\item requisitos, execução, gestão, medição e pagamento\item somente nome do fornecedor\item apenas dotação orçamentária\item senha de sistemas\item marca específica obrigatoriamente\end{enumerate}\end{questao}
\begin{questao}{\nivel{N1} 15.04}SLA significa, em contexto contratual de TI:\begin{enumerate}[label=\Alph*)]\item acordo/nível de serviço com metas e critérios mensuráveis\item sistema de licitação automática\item assinatura eletrônica\item licença de software anual\item lista de ativos\end{enumerate}\end{questao}
\begin{questao}{\nivel{N1} 15.05}Pesquisa de preços tem por objetivo principal:\begin{enumerate}[label=\Alph*)]\item estimar valor de mercado da contratação com base documentada e crítica\item garantir contratação do menor preço isolado\item selecionar fornecedor antecipadamente\item substituir ETP\item eliminar competição\end{enumerate}\end{questao}
\begin{questao}{\nivel{N1} 15.06}Na fiscalização, termo de recebimento definitivo busca registrar:\begin{enumerate}[label=\Alph*)]\item aceitação após verificação das condições previstas\item mera entrega física sem conferência\item orçamento inicial\item pesquisa de mercado\item abertura de chamado\end{enumerate}\end{questao}
\begin{questao}{\nivel{N1} 15.07}Gerenciamento de riscos na contratação deve:\begin{enumerate}[label=\Alph*)]\item acompanhar o ciclo da contratação, não ser exercício único no início\item ocorrer apenas após pagamento\item ser responsabilidade exclusiva da contratada\item substituir fiscalização\item ser dispensado em TIC\end{enumerate}\end{questao}
\begin{questao}{\nivel{N1} 15.08}Lock-in tecnológico significa:\begin{enumerate}[label=\Alph*)]\item dependência que dificulta migração/substituição de fornecedor ou tecnologia\item criptografia de disco\item falha de rede\item bloqueio de usuário\item indisponibilidade temporária\end{enumerate}\end{questao}
\begin{questao}{\nivel{N1} 15.09}Um indicador contratual deve possuir:\begin{enumerate}[label=\Alph*)]\item fórmula, fonte, período, regra de medição e meta\item apenas nome\item apenas cor em dashboard\item somente opinião do fiscal\item nenhum critério de exclusão\end{enumerate}\end{questao}
\begin{questao}{\nivel{N1} 15.10}Segregação de funções busca:\begin{enumerate}[label=\Alph*)]\item reduzir concentração incompatível de decisão, execução e validação em uma mesma pessoa\item aumentar burocracia por si só\item impedir qualquer automação\item eliminar gestores\item substituir auditoria\end{enumerate}\end{questao}

\subsection*{N2 — Aplicação}
\begin{questao}{\nivel{N2} 15.11}Um ETP escolhe produto específico antes de comparar alternativas e depois apenas justifica a escolha. Qual problema?\begin{enumerate}[label=\Alph*)]\item inversão do planejamento e possível direcionamento\item excesso de competição\item falta de SLA exclusivamente\item ausência de backup\item problema de TCP\end{enumerate}\end{questao}
\begin{questao}{\nivel{N2} 15.12}Um contrato paga 100 postos mensais, mas não mede resultados nem demanda efetiva. O principal risco é:\begin{enumerate}[label=\Alph*)]\item pagamento por disponibilidade de mão de obra sem relação clara com resultado/necessidade\item falta de normalização\item ausência de VLAN\item excesso de qualidade\item RPO baixo\end{enumerate}\end{questao}
\begin{questao}{\nivel{N2} 15.13}Disponibilidade contratual de 99,9\% em 30 dias 24x7 permite aproximadamente:\begin{enumerate}[label=\Alph*)]\item 43,2 minutos de indisponibilidade\item 4,32 horas\item 432 minutos\item 14,4 minutos\item zero minuto\end{enumerate}\end{questao}
\begin{questao}{\nivel{N2} 15.14}Uma métrica exclui ``manutenções programadas'', mas o contrato não define limite ou aviso. O risco é:\begin{enumerate}[label=\Alph*)]\item manipulação da disponibilidade por exclusões discricionárias\item falta de criptografia\item ausência de licitação\item excesso de auditoria\item impossibilidade de backup\end{enumerate}\end{questao}
\begin{questao}{\nivel{N2} 15.15}A contratada produz a única planilha que determina seu próprio pagamento, sem evidência independente. Isso fragiliza:\begin{enumerate}[label=\Alph*)]\item confiabilidade e independência da medição\item normalização do banco\item rede física\item assinatura digital\item arquitetura de software\end{enumerate}\end{questao}
\begin{questao}{\nivel{N2} 15.16}Pesquisa de preços com três propostas idênticas de empresas do mesmo grupo econômico deve ser:\begin{enumerate}[label=\Alph*)]\item questionada quanto à independência e representatividade\item aceita automaticamente por ter três cotações\item considerada prova de competição\item usada sem memória de cálculo\item tratada como inexigibilidade\end{enumerate}\end{questao}
\begin{questao}{\nivel{N2} 15.17}Contratações públicas recentes comparáveis foram ignoradas e usaram-se só cotações mais altas. A auditoria deve verificar:\begin{enumerate}[label=\Alph*)]\item justificativa das fontes e aderência aos parâmetros da IN 65\item apenas assinatura do fiscal\item somente quantidade de páginas\item código-fonte do sistema\item versão do Windows\end{enumerate}\end{questao}
\begin{questao}{\nivel{N2} 15.18}Contrato SaaS deve prever plano de saída principalmente para:\begin{enumerate}[label=\Alph*)]\item assegurar portabilidade de dados, continuidade e transição\item aumentar lock-in\item impedir exportação\item eliminar documentação\item dispensar backup\end{enumerate}\end{questao}
\begin{questao}{\nivel{N2} 15.19}Um aceite funcional deve ser baseado em:\begin{enumerate}[label=\Alph*)]\item critérios objetivos e evidências de testes/entrega\item opinião informal da contratada\item pagamento antecipado\item marca do produto\item quantidade de reuniões\end{enumerate}\end{questao}
\begin{questao}{\nivel{N2} 15.20}Se serviço crítico exige RTO de 2 h, o contrato deve conter:\begin{enumerate}[label=\Alph*)]\item responsabilidades, arquitetura/controles e evidência de testes compatíveis com o RTO\item somente preço mensal\item apenas cláusula de confidencialidade\item nenhum procedimento de desastre\item RPO necessariamente igual ao RTO\end{enumerate}\end{questao}
\begin{questao}{\nivel{N2} 15.21}Uma licença anual foi contratada para 5.000 usuários, mas só 1.200 estão ativos. A auditoria deve avaliar:\begin{enumerate}[label=\Alph*)]\item dimensionamento, uso efetivo e oportunidade de otimização\item apenas se a nota foi paga\item somente segurança física\item inexistência de contrato\item apenas marca da licença\end{enumerate}\end{questao}
\begin{questao}{\nivel{N2} 15.22}Uma penalidade contratual deve ser aplicada:\begin{enumerate}[label=\Alph*)]\item segundo previsão contratual/legal, contraditório e evidência do descumprimento\item automaticamente sem processo\item somente se fornecedor concordar\item sem medir o SLA\item por qualquer atraso informal\end{enumerate}\end{questao}
\begin{questao}{\nivel{N2} 15.23}A prorrogação de contrato contínuo deve considerar:\begin{enumerate}[label=\Alph*)]\item vantajosidade, desempenho, condições legais e risco de transição\item apenas conveniência do fornecedor\item ausência de fiscalização\item preço original sem atualização\item obrigação automática de prorrogar\end{enumerate}\end{questao}
\begin{questao}{\nivel{N2} 15.24}Em contrato que trata dados pessoais, deve-se definir:\begin{enumerate}[label=\Alph*)]\item papéis, instruções de tratamento, segurança, incidentes, retenção e destino dos dados\item apenas endereço do datacenter\item somente preço\item dispensa de LGPD\item acesso irrestrito da contratada\end{enumerate}\end{questao}
\begin{questao}{\nivel{N2} 15.25}Uma ordem de serviço deve permitir:\begin{enumerate}[label=\Alph*)]\item rastrear escopo, prazo, responsável, critérios e resultado demandado\item executar qualquer atividade sem autorização\item alterar objeto do contrato livremente\item eliminar aceite\item ocultar custo\end{enumerate}\end{questao}

\subsection*{N3 — Nível de concurso}
\begin{questao}{\nivel{N3} 15.26}Um contrato mede SLA pelo próprio sistema da contratada, mas o órgão recebe logs brutos assinados e recalcula amostras. Isso:\begin{enumerate}[label=\Alph*)]\item aumenta verificabilidade e reduz dependência da medição unilateral\item elimina qualquer risco\item torna auditoria desnecessária\item substitui contrato\item invalida o SLA\end{enumerate}\end{questao}
\begin{questao}{\nivel{N3} 15.27}O fornecedor classifica 80\% dos incidentes como ``força maior'' sem critérios contratuais, excluindo-os do SLA. A falha é:\begin{enumerate}[label=\Alph*)]\item ausência de regra objetiva para exclusões e potencial manipulação do indicador\item falta de RAID\item excesso de logs\item ausência de Scrum\item erro de banco\end{enumerate}\end{questao}
\begin{questao}{\nivel{N3} 15.28}Uma especificação exige processador de marca/modelo exato sem justificativa técnica. O risco é:\begin{enumerate}[label=\Alph*)]\item restrição indevida à competição/direcionamento\item aumento de interoperabilidade automaticamente\item redução de custo garantida\item ausência de objeto\item melhoria da pesquisa de preços\end{enumerate}\end{questao}
\begin{questao}{\nivel{N3} 15.29}Um contrato de desenvolvimento paga por ponto de função, mas entrega código com alto retrabalho. O auditor deve:\begin{enumerate}[label=\Alph*)]\item verificar se métrica de tamanho está acompanhada de critérios de qualidade, aceite e resultado\item concluir que ponto de função mede qualidade\item ignorar defeitos\item pagar apenas por linhas de código\item eliminar testes\end{enumerate}\end{questao}
\begin{questao}{\nivel{N3} 15.30}Em contratação de cloud, estimar apenas mensalidade-base e ignorar egress, suporte e crescimento viola principalmente:\begin{enumerate}[label=\Alph*)]\item análise de custo total e planejamento da solução\item sigilo\item integridade referencial\item disponibilidade do edital\item processo legislativo\end{enumerate}\end{questao}
\begin{questao}{\nivel{N3} 15.31}Um contrato exige 99,99\% de disponibilidade, mas não há arquitetura redundante nem orçamento compatível. Isso indica:\begin{enumerate}[label=\Alph*)]\item requisito possivelmente inexequível ou desconectado do desenho/custo\item SLA forte sempre reduz custo\item disponibilidade independe de arquitetura\item contrato é automaticamente nulo\item RPO igual a zero\end{enumerate}\end{questao}
\begin{questao}{\nivel{N3} 15.32}O fiscal atesta serviço com base em e-mail ``tudo certo'', sem verificar chamados e logs. A deficiência é:\begin{enumerate}[label=\Alph*)]\item aceite sem evidência suficiente e objetiva\item excesso de formalismo\item boa segregação\item redundância de controle\item falta de nuvem\end{enumerate}\end{questao}
\begin{questao}{\nivel{N3} 15.33}A mesma pessoa elabora requisitos, conduz seleção e atesta recebimento sem revisão. O risco é:\begin{enumerate}[label=\Alph*)]\item concentração de funções e conflito de interesses\item perda de disponibilidade apenas\item falta de teste unitário\item excesso de competição\item integridade criptográfica\end{enumerate}\end{questao}
\begin{questao}{\nivel{N3} 15.34}O Mapa de Gerenciamento de Riscos não é atualizado após mudança tecnológica relevante. O problema é:\begin{enumerate}[label=\Alph*)]\item gestão de risco tratada como documento estático, sem refletir risco atual\item excesso de revisão\item risco eliminado\item contrato automaticamente encerrado\item problema de DNS\end{enumerate}\end{questao}
\begin{questao}{\nivel{N3} 15.35}Um fornecedor retém código-fonte e documentação essenciais sem cláusula clara de propriedade/licenciamento. O risco é:\begin{enumerate}[label=\Alph*)]\item dependência e dificuldade de continuidade/manutenção\item aumento de portabilidade\item redução de lock-in\item ausência de dados\item apenas problema de SLA\end{enumerate}\end{questao}
\begin{questao}{\nivel{N3} 15.36}Pesquisa de preços mistura assinatura SaaS mensal com licença perpétua sem normalizar horizonte e escopo. A consequência é:\begin{enumerate}[label=\Alph*)]\item comparação inválida por bases econômicas distintas\item maior precisão\item mediana sempre corrige\item preço deixa de importar\item inexigibilidade automática\end{enumerate}\end{questao}
\begin{questao}{\nivel{N3} 15.37}Um indicador mede ``chamados resolvidos'', mas a contratada fecha e reabre chamados para melhorar resultado. O controle adequado é:\begin{enumerate}[label=\Alph*)]\item incluir reabertura, aceite do usuário e regras anti-gaming\item aumentar meta de fechamentos\item remover logs\item aceitar dados da contratada\item trocar SLA por horas\end{enumerate}\end{questao}
\begin{questao}{\nivel{N3} 15.38}Na transição contratual, contas privilegiadas do fornecedor antigo permanecem ativas. O risco é:\begin{enumerate}[label=\Alph*)]\item acesso indevido após término e falha de offboarding\item indisponibilidade do novo fornecedor apenas\item ausência de pesquisa de preço\item erro de licitação\item melhoria da continuidade\end{enumerate}\end{questao}
\begin{questao}{\nivel{N3} 15.39}Contrato de backup prevê retenção, mas não exige restore test. A avaliação correta é:\begin{enumerate}[label=\Alph*)]\item retenção não prova recuperabilidade; testes devem validar restauração\item retenção torna teste desnecessário\item backup é igual a RAID\item RTO não importa\item RPO é preço\end{enumerate}\end{questao}
\begin{questao}{\nivel{N3} 15.40}Uma contratação emergencial é renovada repetidamente por falta de planejamento. A auditoria deve investigar:\begin{enumerate}[label=\Alph*)]\item causas da emergência recorrente e falha de planejamento/gestão\item apenas preço atual\item somente segurança física\item linguagem de programação\item número de funcionários do fornecedor\end{enumerate}\end{questao}

\subsection*{N4 — Avançado / Auditoria}
\begin{questao}{\nivel{N4} 15.41}Plataforma SaaS crítica foi contratada sem exportação padronizada nem documentação de APIs. Após dois anos, migração é estimada em 18 meses. O risco principal é:\begin{enumerate}[label=\Alph*)]\item lock-in tecnológico/contratual e ausência de estratégia de saída\item fragmentação IP\item colisão de hash\item deadlock\item falta de compressão\end{enumerate}\end{questao}
\begin{questao}{\nivel{N4} 15.42}Contrato de sustentação paga valor fixo, não possui catálogo de serviços, SLA é genérico e não há penalidades registradas por 24 meses. O auditor deve prioritariamente:\begin{enumerate}[label=\Alph*)]\item testar demanda, medição, qualidade, cumprimento e relação entre pagamento e resultado\item concluir que não houve falha porque não houve penalidade\item verificar apenas assinatura\item aceitar relatórios da contratada\item ignorar chamados\end{enumerate}\end{questao}
\begin{questao}{\nivel{N4} 15.43}A contratada administra o sistema de monitoramento e pode alterar históricos antes da medição mensal. O controle mais adequado é:\begin{enumerate}[label=\Alph*)]\item logs/telemetria imutáveis ou sob controle do órgão, com reconciliação independente\item confiar na boa-fé\item aumentar SLA\item remover monitoramento\item pagar antecipado\end{enumerate}\end{questao}
\begin{questao}{\nivel{N4} 15.44}Uma licitação de SOC exige equipe 24x7, mas não define casos de uso, cobertura de logs, MTTD/MTTR ou qualidade de investigação. O risco é:\begin{enumerate}[label=\Alph*)]\item contratar presença de pessoas sem resultado de segurança mensurável\item excesso de requisitos\item ausência de firewall apenas\item falta de licitação\item custo necessariamente baixo\end{enumerate}\end{questao}
\begin{questao}{\nivel{N4} 15.45}O órgão contrata 1 PB de storage ``para crescimento futuro'', sem histórico de consumo. Após 18 meses, usa 180 TB. O achado deve focar:\begin{enumerate}[label=\Alph*)]\item dimensionamento sem premissas, possível supercontratação e custo de ociosidade\item falta de RAID exclusivamente\item baixa segurança física\item inexistência de software\item uso insuficiente de nuvem\end{enumerate}\end{questao}
\begin{questao}{\nivel{N4} 15.46}Um reequilíbrio é pedido por aumento previsível de custos que já integravam risco normal do negócio. O auditor deve:\begin{enumerate}[label=\Alph*)]\item examinar fundamento jurídico, matriz de risco e demonstração do evento, sem presumir direito automático\item aceitar qualquer pedido\item negar qualquer pedido sempre\item usar apenas inflação geral\item ignorar contrato\end{enumerate}\end{questao}
\begin{questao}{\nivel{N4} 15.47}Contrato de IA permite ao fornecedor usar dados do órgão para ``melhoria de modelos'' sem limites. Qual risco?\begin{enumerate}[label=\Alph*)]\item tratamento secundário, confidencialidade, LGPD, propriedade e vazamento de informação\item apenas custo computacional\item somente disponibilidade\item nenhuma, por ser IA\item problema de VLAN\end{enumerate}\end{questao}
\begin{questao}{\nivel{N4} 15.48}Após término de contrato cloud, o fornecedor afirma ter apagado dados, mas não há evidência de eliminação de backups. O controle esperado é:\begin{enumerate}[label=\Alph*)]\item procedimento contratual de retorno/eliminação, prazos, evidência/certificação e tratamento de cópias\item confiar verbalmente\item deixar dados indefinidamente\item apagar apenas contas de usuário\item remover logs do órgão\end{enumerate}\end{questao}
\begin{questao}{\nivel{N4} 15.49}Um SLA de atendimento mede resposta inicial, mas não solução; a contratada responde em 5 min e deixa incidentes abertos por dias. A correção é:\begin{enumerate}[label=\Alph*)]\item combinar métricas de resposta, restauração/resolução, severidade e qualidade\item reduzir resposta para 1 min apenas\item eliminar SLA\item pagar por chamados abertos\item medir somente e-mail\end{enumerate}\end{questao}
\begin{questao}{\nivel{N4} 15.50}O gestor deseja aceitar entrega com defeitos críticos para não perder saldo orçamentário. A atuação adequada do fiscal é:\begin{enumerate}[label=\Alph*)]\item aplicar critérios de aceite e registrar não conformidade; conveniência orçamentária não substitui evidência de entrega\item aceitar e corrigir depois sem registro\item alterar critério retroativamente\item deixar fornecedor atestar\item pagar integralmente sem teste\end{enumerate}\end{questao}


\section{Treino discursivo}

\begin{discursiva}[fiscalização contratual]
Em auditoria de contrato de sustentação, verifica-se pagamento mensal fixo, ausência de catálogo de serviços, SLAs genéricos, chamados encerrados pela contratada sem aceite do usuário e nenhuma penalidade nos últimos 24 meses. Em até 30 linhas, apresente riscos, procedimentos de auditoria e recomendações.
\end{discursiva}

\begin{discursiva}[lock-in e nuvem]
Um órgão pretende renovar plataforma SaaS porque a migração é considerada ``impossível''. Não há exportação testada, documentação de APIs nem plano de transição. Em até 30 linhas, estruture procedimentos de auditoria para avaliar lock-in, economicidade da renovação e medidas de saída.
\end{discursiva}

\section{Gabarito e resoluções comentadas — Fiscalização de Contratos de TI}

\begin{solucao}{15.01}\textbf{Gabarito: A.} A IN SGD/ME nº 94/2022 organiza o macroprocesso em Planejamento da Contratação, Seleção do Fornecedor e Gestão do Contrato. ETP e TR são artefatos/etapas dentro do planejamento, não fases equivalentes ao macroprocesso.\end{solucao}
\begin{solucao}{15.02}\textbf{Gabarito: A.} O ETP parte da necessidade, compara alternativas e sustenta a viabilidade da solução. Não substitui contrato, recebimento ou aplicação de sanções.\end{solucao}
\begin{solucao}{15.03}\textbf{Gabarito: A.} O TR transforma a solução escolhida em requisitos executáveis e mensuráveis, incluindo modelo de execução/gestão e pagamento. Marca específica sem justificativa pode restringir competição.\end{solucao}
\begin{solucao}{15.04}\textbf{Gabarito: A.} SLA define níveis de serviço com métricas, metas, janelas e consequências conforme contrato. Não é sistema de licitação nem licença.\end{solucao}
\begin{solucao}{15.05}\textbf{Gabarito: A.} Pesquisa de preços sustenta estimativa razoável, documentando fontes, comparabilidade e método. Não escolhe fornecedor antecipadamente nem se resume ao menor número encontrado.\end{solucao}
\begin{solucao}{15.06}\textbf{Gabarito: A.} Recebimento definitivo deve decorrer de verificação de conformidade. Mera entrega física pode fundamentar recebimento provisório, conforme objeto/procedimento, mas não substitui aceite final.\end{solucao}
\begin{solucao}{15.07}\textbf{Gabarito: A.} Riscos mudam com solução, seleção, execução e mudança contratual; o mapa precisa ser vivo. Tratar risco só no início reduz valor do controle.\end{solucao}
\begin{solucao}{15.08}\textbf{Gabarito: A.} Lock-in é dependência de fornecedor, formato, API, licença ou tecnologia que torna saída cara/demorada. Não é bloqueio de conta nem incidente de rede.\end{solucao}
\begin{solucao}{15.09}\textbf{Gabarito: A.} Indicador auditável precisa ser reproduzível: fórmula, fonte, período, meta e regras de inclusão/exclusão. Opinião do fiscal sem critério é insuficiente.\end{solucao}
\begin{solucao}{15.10}\textbf{Gabarito: A.} Segregação reduz possibilidade de uma pessoa definir, executar e validar operação sem revisão. Não é burocracia como fim em si.\end{solucao}
\begin{solucao}{15.11}\textbf{Gabarito: A.} ETP deve avaliar alternativas antes de concluir pela solução. Decisão prévia seguida de justificativa pode mascarar direcionamento e inviabilizar análise de vantajosidade.\end{solucao}
\begin{solucao}{15.12}\textbf{Gabarito: A.} Medir apenas postos pode remunerar ociosidade e deslocar risco de produtividade ao órgão. O modelo deve relacionar pagamento a entregas/resultados quando tecnicamente possível.\end{solucao}
\begin{solucao}{15.13}\textbf{Gabarito: A.} Em 30 dias há 43.200 minutos. A indisponibilidade para 99,9\% é $0,1\%\times43.200=43,2$ min. A regra real depende da janela e exclusões definidas.\end{solucao}
\begin{solucao}{15.14}\textbf{Gabarito: A.} Sem definição de manutenção programada, o fornecedor pode excluir períodos relevantes e inflar a disponibilidade. Métrica precisa de regra objetiva anterior à medição.\end{solucao}
\begin{solucao}{15.15}\textbf{Gabarito: A.} A parte economicamente interessada controla a única evidência, o que fragiliza independência e verificabilidade. O órgão deve ter logs, acesso ou mecanismo de reconciliação.\end{solucao}
\begin{solucao}{15.16}\textbf{Gabarito: A.} Três propostas não garantem mercado representativo quando há vínculo entre cotantes. O auditor deve verificar independência, comparabilidade e possíveis sinais de conluio.\end{solucao}
\begin{solucao}{15.17}\textbf{Gabarito: A.} A IN 65 estabelece parâmetros e prioriza fontes públicas relevantes quando disponíveis. Ignorá-las sem justificativa pode distorcer a estimativa.\end{solucao}
\begin{solucao}{15.18}\textbf{Gabarito: A.} Plano de saída trata exportação, formatos, documentação, prazos, suporte à transição, apagamento e continuidade, reduzindo lock-in.\end{solucao}
\begin{solucao}{15.19}\textbf{Gabarito: A.} Aceite deve confrontar entrega com critérios previamente definidos. A declaração da própria contratada não é suficiente para provar conformidade.\end{solucao}
\begin{solucao}{15.20}\textbf{Gabarito: A.} RTO precisa ser traduzido em arquitetura, responsabilidades, procedimentos e testes. Cláusula abstrata sem capacidade demonstrável é frágil. RPO é objetivo distinto.\end{solucao}
\begin{solucao}{15.21}\textbf{Gabarito: A.} Uso de 24\% das licenças indica necessidade de revisar premissas, sazonalidade, crescimento e alternativas de licenciamento. O pagamento formal não prova economicidade.\end{solucao}
\begin{solucao}{15.22}\textbf{Gabarito: A.} Penalidade exige previsão, fato comprovado, nexo com obrigação e devido processo/contraditório aplicável. Automatismo sem evidência é inadequado.\end{solucao}
\begin{solucao}{15.23}\textbf{Gabarito: A.} Prorrogação deve ser motivada e demonstrar aderência às condições legais e vantajosidade, considerando desempenho e risco de transição. Não é direito automático.\end{solucao}
\begin{solucao}{15.24}\textbf{Gabarito: A.} Contrato deve disciplinar finalidade/instruções, acesso, segurança, incidentes, suboperadores, retenção, transferência e eliminação conforme papéis e bases aplicáveis.\end{solucao}
\begin{solucao}{15.25}\textbf{Gabarito: A.} OS/OFB precisa permitir rastrear demanda e verificar resultado. Autorizações vagas aumentam risco de escopo não controlado e pagamento indevido.\end{solucao}
\begin{solucao}{15.26}\textbf{Gabarito: A.} Receber logs brutos protegidos e recalcular amostras melhora independência da evidência. Isso não elimina todo risco, pois configuração e completude dos logs também precisam ser validadas.\end{solucao}
\begin{solucao}{15.27}\textbf{Gabarito: A.} Exceções sem taxonomia e evidência podem transformar qualquer falha em exclusão. O contrato deve definir força maior, manutenção, terceiros e processo de validação.\end{solucao}
\begin{solucao}{15.28}\textbf{Gabarito: A.} Exigir marca/modelo sem fundamento funcional/técnico pode restringir competição. O requisito deve descrever desempenho e compatibilidade ou justificar exceção.\end{solucao}
\begin{solucao}{15.29}\textbf{Gabarito: A.} Ponto de função mede tamanho funcional, não qualidade do código. Aceite deve combinar tamanho/entrega com defeitos, testes, segurança e critérios de qualidade.\end{solucao}
\begin{solucao}{15.30}\textbf{Gabarito: A.} Cloud tem componentes variáveis de custo. Ignorar egress, suporte, crescimento, logs e serviços associados subestima TCO e prejudica comparação de alternativas.\end{solucao}
\begin{solucao}{15.31}\textbf{Gabarito: A.} SLO extremo exige redundância e custo compatíveis. Meta arbitrária pode elevar preço ou ser inexequível. Planejamento deve relacionar necessidade, impacto e capacidade técnica.\end{solucao}
\begin{solucao}{15.32}\textbf{Gabarito: A.} E-mail genérico não demonstra quantidade, qualidade ou cumprimento de SLA. O fiscal precisa de evidência suficiente e apropriada ao objeto.\end{solucao}
\begin{solucao}{15.33}\textbf{Gabarito: A.} Concentrar especificação, seleção e aceite reduz checks and balances e aumenta risco de favorecimento/erro não detectado.\end{solucao}
\begin{solucao}{15.34}\textbf{Gabarito: A.} Risco é dinâmico. Mudança de arquitetura, fornecedor, legislação ou ameaça deve provocar reavaliação e atualização de respostas.\end{solucao}
\begin{solucao}{15.35}\textbf{Gabarito: A.} Código/documentação sem direitos e condições claras cria dependência do fornecedor e pode impedir manutenção por terceiros. O planejamento deve tratar propriedade/licenciamento e transição.\end{solucao}
\begin{solucao}{15.36}\textbf{Gabarito: A.} Mensalidade e licença perpétua têm horizontes e componentes diferentes; devem ser normalizadas por período, escopo, suporte e TCO. Média/mediana não corrige dados incomparáveis.\end{solucao}
\begin{solucao}{15.37}\textbf{Gabarito: A.} Métrica unilateral de fechamento pode ser manipulada. Reabertura, tempo até solução e aceite do usuário reduzem gaming e alinham indicador ao resultado.\end{solucao}
\begin{solucao}{15.38}\textbf{Gabarito: A.} Offboarding incompleto mantém acesso de terceiro sem necessidade. Deve haver revogação de contas, tokens, VPNs, certificados e privilégios no encerramento/transição.\end{solucao}
\begin{solucao}{15.39}\textbf{Gabarito: A.} Retenção demonstra existência potencial de cópias; só restore test evidencia que cadeia, chaves, catálogo e procedimentos realmente recuperam o serviço.\end{solucao}
\begin{solucao}{15.40}\textbf{Gabarito: A.} Emergência recorrente pode ser consequência de planejamento inadequado, dependência contratual ou falha de governança. O auditor deve examinar causa, não apenas preço da última contratação.\end{solucao}
\begin{solucao}{15.41}\textbf{Gabarito: A.} Sem portabilidade e APIs documentadas, o órgão perde poder de saída. Plano de exit deve estar no ETP/TR antes da contratação, não apenas quando o fornecedor aumenta custos.\end{solucao}
\begin{solucao}{15.42}\textbf{Gabarito: A.} Ausência de catálogo e SLA mensurável impede demonstrar o que foi comprado e entregue. O auditor deve cruzar demanda, pagamentos, chamados, equipe e resultados.\end{solucao}
\begin{solucao}{15.43}\textbf{Gabarito: A.} Evidência usada para pagar não deve poder ser reescrita unilateralmente. Telemetria protegida/imutável, sob domínio do órgão ou reconciliada, aumenta confiabilidade.\end{solucao}
\begin{solucao}{15.44}\textbf{Gabarito: A.} SOC deve ser contratado por capacidade e resultado de detecção/resposta, com cobertura de fontes, casos de uso, tempos, qualidade e evidências. Apenas quantidade de pessoas não mede segurança.\end{solucao}
\begin{solucao}{15.45}\textbf{Gabarito: A.} Ociosidade de cerca de 82\% sugere dimensionamento sem base ou antecipação excessiva. O achado deve examinar premissas, custo evitável e flexibilidade contratual.\end{solucao}
\begin{solucao}{15.46}\textbf{Gabarito: A.} Reequilíbrio depende do regime legal, matriz de risco e comprovação de evento/impacto. Risco ordinário do contratado não gera direito automático.\end{solucao}
\begin{solucao}{15.47}\textbf{Gabarito: A.} Uso secundário de dados pode violar finalidade, confidencialidade e direitos de propriedade, além de criar risco de memorização/exposição. Deve haver limites contratuais claros.\end{solucao}
\begin{solucao}{15.48}\textbf{Gabarito: A.} Encerramento precisa prever devolução/exportação, eliminação de produção e backups conforme retenção, evidência da destruição e exceções legais.\end{solucao}
\begin{solucao}{15.49}\textbf{Gabarito: A.} Resposta inicial pode ser cosmeticamente rápida enquanto usuário permanece indisponível. Métricas devem refletir restauração/resolução e severidade.\end{solucao}
\begin{solucao}{15.50}\textbf{Gabarito: A.} Fiscal não deve alterar critério para consumir orçamento. Pagamento deve corresponder ao objeto aceito; não conformidade precisa ser registrada e tratada conforme contrato.\end{solucao}

